\scanpage{59}
\mspage{10}
Passons au cas de $Q_i$ -- la définition \unsure{évidente} \deleted{\unsure{géométrique}} des
\[
\psi_r=\psi_i^{\omega,\varepsilon}:\quad\Pi_{0,3}-\Pi_i\deleted{(\unsure{Q_i})}
\]
est
\[
\inserted{\psi_r=}\psi_i^{\omega,\varepsilon}=\Pi_1(\underbrace{u_{i,\omega,\varepsilon}}_{u_r};Q_\infty)
\]
\margin{\uline{NB} on identifie la \inserted{source $\Pi_1(U_{0,3};Q_\infty)$} \deleted{\unsure{source}} à $\Pi_{03}\stackrel{\mathrm{def}}=\Pi_1(U_{0,3};P_+)$ à l'aide du chemin standard $d_\infty^+$ de $P_+$ à $Q_\infty$.}
\unsure{Donc} $u_r=u_{i,\omega,\varepsilon}$ est l'unique isom de $U_{0,3}$, \inserted{dans $\Sigma^*$} qui transforme le repère \deleted{fixé} standard $(\infty,\omega_0=+1,\unsure{\bar\omega|\varpi}_0=+1\text{ donc }\varepsilon=\omega_0\wedge\unsure{\bar\omega|\varpi}_0=+1)$ en $r$ i.e. qui transforme $P_+$ en $P_\varepsilon$, $Q_0$ en $Q_i$, $Q_1$ en $Q_{\omega i}$. \deleted{On} Les épinglages \deleted{\unsure{par un repère fixé}} de $\Pi_i$ \inserted{$\psi_i^{\omega,\varepsilon}$ ($i$ fixé, $\omega,\varepsilon$ variables)} correspondent aux $u_r$ qui transforment $Q_0$ en $Q_i$, ils sont donc déterminés par composition avec des éléments
\[
g_0\in\Gamma_0=\mathcal G_3\times\{1,\tau\}
\]
qui stabilisent $Q_0$ -- ces éléments \deleted{engendrent} font un groupe $\simeq\mathbb Z/2\times\mathbb Z/2$ engendré par $\sigma_\infty,\tau$, \deleted{et} On trouve encore
\[
\psi_{r'}^{\deleted{\unsure{\omega}}}=\psi_r\circ\deleted{\unsure{\psi}}\Pi_1(g_0;Q_\infty)\quad\text{si }r'=r\circ g_0
\]
où ici $\psi_{r'},\psi_r$ désignent les iso. d'épinglage relatifs à $Q_i$. Il suffit d'expliciter ces formules dans les cas de $g_0=\sigma_\infty,g_0=\tau$, ce qui revient à calculer leur effet sur $\deleted{\unsure{\Pi_{0,3}=}}\Pi_1\deleted{\unsure{\varepsilon}}(U_{0,3};\unsure{Q_\infty|R_\infty})$, quand on identifie celui-ci à $\Pi_{03}=\Pi_1(U_{0,3};P_+)$, grâce au chemin $d_\infty^+$ de $P_+$ à $Q_\infty$. On sait que l'on peut
\margin{58}
